\section{Páginas complementarias de las diapositivas}

\subsection{Dimensiones de alineación de comportamiento}
\clearpage
\noindent\textbf{Dimensiones de alineación de comportamiento}\\
Esta diapositiva describe las relaciones de dependencia entre la estructura de tres niveles de recompensa explícita, preferencia implícita y guardrail de seguridad, lo cual corresponde al reward signal breakdown de este texto.
\begin{center}
\includegraphics[width=0.92\textwidth,page=11]{08_cs224r_reward_learning_2025.pdf}
\end{center}

\subsection{Proceso de Agent Swarm}
\clearpage
\noindent\textbf{Proceso de Agent Swarm}\\
Esta figura muestra cómo varios Agent comparten el reward trace y cómo se dividen el trabajo para procesar la metric evaluation, lo cual se corresponde con el Operations Playbook de este capítulo.
\begin{center}
\includegraphics[width=0.92\textwidth,page=12]{08_cs224r_reward_learning_2025.pdf}
\end{center}

\subsection{Control de calidad de datos}
\clearpage
\noindent\textbf{Control de calidad de datos}\\
La estructura por niveles de Data curation pipeline, filter pipes y human annotation loops puede servir de referencia para la sección de «datos por lotes» de este capítulo.
\begin{center}
\includegraphics[width=0.92\textwidth,page=13]{08_cs224r_reward_learning_2025.pdf}
\end{center}

\subsection{Mecanismos de gobernanza}
\clearpage
\noindent\textbf{Mecanismos de gobernanza}\\
Explica la conexión entre audit logs, human review y policy rollback, lo cual corresponde al checklist del capítulo de gobernanza.
\begin{center}
\includegraphics[width=0.92\textwidth,page=14]{08_cs224r_reward_learning_2025.pdf}
\end{center}

\subsection{Restricciones de Sim-to-Real}
\clearpage
\noindent\textbf{Restricciones de Sim-to-Real}\\
Las Slides enfatizan la combinación de domain randomization y reward smoothing, contenido que agregamos en la sección de «Sim-to-Real».
\begin{center}
\includegraphics[width=0.92\textwidth,page=15]{08_cs224r_reward_learning_2025.pdf}
\end{center}

\subsection{Ops Playbook}
\clearpage
\noindent\textbf{Ops Playbook}\\
\begin{center}
\includegraphics[width=0.92\textwidth,page=16]{08_cs224r_reward_learning_2025.pdf}
\end{center}
\noindent La figura muestra los distintos checkpoint y runbook, lo cual respalda directamente el checklist propuesto en el «Observability Playbook» de este capítulo.

\subsection{Evidence Matrix}
\clearpage
\noindent\textbf{Evidence Matrix}\\
Enumera el mapeo entre reward signal, policy y safety eval, lo cual facilita evaluar la evidence quality de distintas entradas.
\begin{center}
\includegraphics[width=0.92\textwidth,page=17]{08_cs224r_reward_learning_2025.pdf}
\end{center}

\subsection{Direcciones futuras de investigación}
\clearpage
\noindent\textbf{Direcciones futuras de investigación}\\
Enumera direcciones por explorar como CAI, Verifier y Adaptive Reward, una extensión directa de la sección «direcciones futuras» de esta clase.
\begin{center}
\includegraphics[width=0.92\textwidth,page=18]{08_cs224r_reward_learning_2025.pdf}
\end{center}
